% ---
% id: receiver/prop:res5
% kind: prop
% title: "Properties of the Hermiticity-and-positivity defect"
% summary: >-
%   The defect of a matrix vanishes exactly when it is Hermitian and positive
%   semidefinite, and bounds how far its expectation values stray below zero or
%   off the real axis. Every emitted block before the dark counts has defect
%   zero exactly, so the fifth residual measures the evaluation alone.
% prominence: standard
% status: implemented
% uses: [def:residuals, thm:coarse-graining, def:setting-family]
% assumes: []
% refs: []
% provenance: newly-derived
% implements: certiqs_key.receiver.numerics:hermiticity_positivity_defect
% lean:
% ---

\begin{proposition}\label{prop:res5}
Let $X$ be a complex $D\times D$ matrix and $\varrho(X)$ its
\tref{Hermiticity-and-positivity defect} \eqref{eq:rhodef}.
\begin{enumerate}[label=(\roman*),leftmargin=3.0em]
\item $\varrho(X)=0$ if and only if $X=X^\dagger$ and $X\succeq0$.
\item If $X=X^\dagger$ then $\varrho(X)=\big(-\lmin(X)\big)_+$, which is the
  form \eqref{eq:res2} takes on a Gram.
\item For every unit vector $\ket{\psi}$,
  $\Re\bra{\psi}X\ket{\psi}\ge-\varrho(X)$ and
  $\big|\Im\bra{\psi}X\ket{\psi}\big|\le\tfrac12\varrho(X)$. Hence
  $\varrho(X)\le\varepsilon$ bounds every diagonal expectation of $X$ away from
  the negative reals by $\varepsilon$ and off the real axis by
  $\varepsilon/2$.
\item Every block $\blk{M}{\eqn}_{v|t}$ of the coarse-grained elements, before
  any dark-count channel, that the artifact emits satisfies $\varrho=0$
  \emph{exactly}: for
  each $v$, $n$ and $t$, $\blk{M}{\eqn}_{v|t}$ is a diagonal block of the
  positive operator $F_{v|t}$ (Definition~\ref{def:setting-family}\,(a) and
  Theorem~\ref{thm:coarse-graining}), and a diagonal block --- the compression $\Pi X\Pi$ of a positive
  operator $X$ by an orthogonal projector $\Pi$, read on the range of $\Pi$ ---
  is Hermitian and positive semidefinite. So $r_5$ is a residual in exactly the
  sense of the other five: the exact value is $0$, and what is emitted is the
  floating-point evaluation.
\end{enumerate}
\end{proposition}

\begin{proof}
Write $H\coloneqq\tfrac12(X+X^\dagger)$ and $S\coloneqq X-X^\dagger$, so that
$X=H+\tfrac12S$ with $H$ Hermitian and $S$ skew-Hermitian.

(i) If $\varrho(X)=0$ then $\lVert S\rVert=0$, so $X=X^\dagger$ and $H=X$; and
$(-\lmin(X))_+=0$ gives $\lmin(X)\ge0$, i.e.\ $X\succeq0$. Conversely if
$X=X^\dagger\succeq0$ then $S=0$ and $\lmin(X)\ge0$, so both entries of the
maximum vanish.

(ii) $S=0$ gives $\lVert S\rVert=0$ and $H=X$, so the maximum is its second
entry.

(iii) $\bra{\psi}X\ket{\psi}=\bra{\psi}H\ket{\psi}
+\tfrac12\bra{\psi}S\ket{\psi}$. The first term is real, with
$\bra{\psi}H\ket{\psi}\ge\lmin(H)\ge-\varrho(X)$ by the definition of
$\varrho$; the second is purely imaginary, because $iS$ is Hermitian and
therefore $\bra\psi S\ket\psi\in i\mathbb{R}$, with modulus at most
$\tfrac12\lVert S\rVert\le\tfrac12\varrho(X)$. Splitting into real and
imaginary parts gives both inequalities.

(iv) A POVM element is positive semidefinite, and so is any sum of them, so
$F_{v|t}=\sum_{S\in\cg^{-1}(v)}M_{S|t}\succeq0$. Each $F_{v|t}$ is block
diagonal in the photon number (Definition~\ref{def:setting-family}\,(a) and
Theorem~\ref{thm:coarse-graining}), so $\blk{M}{\eqn}_{v|t}$ is the compression
$\Pi_nF_{v|t}\Pi_n$ read on the range of $\Pi_n$. For any $Y\succeq0$ and any
orthogonal projector $\Pi$, $\Pi Y\Pi$ is Hermitian and
$\bra{\psi}\Pi Y\Pi\ket{\psi}=\bra{\Pi\psi}Y\ket{\Pi\psi}\ge0$, so the
compression is Hermitian positive semidefinite. Now apply (i).
\end{proof}
